\chapter{Interacting Conditional-Prediction Agents}
\label{app:interacting-memory}

The conditional model of \cref{app:conditional-memory} becomes an interacting system once a transmitted result has an actual record law. This appendix supplies one such law. Two agents share a normalized preparation and a communication protocol, receive separate private records, and exchange fresh samples from their current result predictions. The receiver conditions on the sender's sampling mechanism, including its inaccessible private evidence. For every finite history outside the usual conditioning null sets, the resulting posterior can be reconstructed from the joint posterior of the private records and one fixed conditional kernel. This is an exact representation before any moment projection.

The common preparation, its cross-agent coupling, the private sensors, and reliable registration of the transcript are constitutive assumptions. Sharing a local bank does not supply them. The distinction between common information and privately evaluated prescriptions is familiar from decentralized control \citep{Nayyar2013CommonInformation}. Here the prescription is fixed to posterior-predictive transmission; the results below concern its normalized law and Bayesian conditioning, without an optimal-control assertion. All variables are defined at one context. No connection, continuous-time flow, or identification with the reciprocal law potential is assumed.

\section{Preparation, information, and transmission}

For $i\in\{0,1\}$, let $\mathsf A_i,\mathsf X_i,\mathsf Y_i$, and $\mathsf D_i$ be standard Borel spaces of model presentations, cues, results, and private records. Write $Z_i=(A_i,X_i,Y_i)$ and $Z=(Z_0,Z_1)$. A proper presentation law $r_i$, a normalized cue kernel $\mu_{i,a}$, and a normalized result kernel $K_{i,a}$ define the local bank
\begin{equation}
 \pi_i(da,du,dv)=r_i(da)\mu_{i,a}(du)K_{i,a}(u,dv).
 \label{imem:bank}
\end{equation}
Choose a probability coupling $P_0$ of $\pi_0$ and $\pi_1$. For example, any coupling $\Xi$ of the two presentation/cue laws gives
\begin{equation}
 P_0(dz_0,dz_1)=\Xi(da_0,du_0,da_1,du_1)
 K_{0,a_0}(u_0,dv_0)K_{1,a_1}(u_1,dv_1).
 \label{imem:coupled-bank}
\end{equation}
Each result factor integrates to one, so this is normalized and has the required local marginals. More general couplings with those marginals are allowed in the general construction.

Let $L_i:\mathsf X_i\rightsquigarrow\mathsf D_i$ be normalized private-record kernels, and put $D=(D_0,D_1)$ and $E=(Z,D)$. The preparation is
\begin{equation}
 \Lambda_0(dz,dd)=P_0(dz_0,dz_1)L_0(u_0,dd_0)L_1(u_1,dd_1).
 \label{imem:preparation}
\end{equation}
Thus the private records are conditionally independent given $Z$. All integrations in \cref{imem:bank,imem:coupled-bank,imem:preparation} are kernel integrations; no Lebesgue density or Gaussian support is required. Agent $i$ observes its value $d_i$, knows the preparation and protocol, and initially has the correlated recognition law
\begin{equation}
 \zeta_{i,0}(da_i,du_i,dv_i)
 =P(Z_i\in\cdot\mid D_i=d_i)
 =s_{i,0}(da_i)q_{i,0}(du_i,dv_i\mid a_i).
 \label{imem:initial-recognition}
\end{equation}
Marginalizing the other agent before conditioning gives exactly the local bank posterior. Presentation/content independence is not imposed.

The preparation may be realized by concealed internal values $Z$ followed by their records. Alternatively, its record marginal can be sampled directly after integrating out $Z$. These realizations have the same communication law below; statements about represented variables refer to the augmented model \cref{imem:preparation}. Normalization does not require an external ontology for those variables.

Let $\mathsf H$ be a standard Borel message space, let $j_t$ be a fixed known alternating sequence of speakers, and let
\begin{equation}
 \mathcal H_t=\sigma(H_1,\ldots,H_t),\qquad
 \mathcal I_{i,t}=\sigma(D_i,H_1,\ldots,H_t).
 \label{imem:information}
\end{equation}
The common and private posteriors, and the local recognition, are
\begin{equation}
 \Gamma_t=P(E\in\cdot\mid\mathcal H_t),\qquad
 \Pi_{i,t}=P(E\in\cdot\mid\mathcal I_{i,t}),\qquad
 \zeta_{i,t}=(Z_i)_\#\Pi_{i,t}.
 \label{imem:posteriors}
\end{equation}
The common posterior is determined from the common transcript; it does not disclose either private value. Regular conditional probabilities exist on these spaces \citep{Kallenberg2021}. Choose their measurable versions in the history and private-record variables, with fixed normalized extensions on conditioning null sets. Each finite history is constructed before its conditional kernels are used to generate the next message. An autonomous map on a larger family of law-valued states additionally requires a jointly measurable choice of disintegration on that family. Finite tables and the dominated formulas below provide it directly.

\constructionheading{Posterior-sampled communication}{imem:protocol}
Fix normalized encoders $C_j:\mathsf Y_j\rightsquigarrow\mathsf H$. At the next step $j=j_{t+1}$, let
\begin{equation}
 S_{j,t}(d_j,dy)=P(Y_j\in dy\mid D_j=d_j,\mathcal H_t),\qquad
 \kappa_{j,t}(d_j,dh)=\int C_j(y,dh)S_{j,t}(d_j,dy).
 \label{imem:message-kernel}
\end{equation}
The speaker draws a fresh $\widetilde Y_{t+1}$ with law $S_{j,t}(D_j,\cdot)$ and encodes it using $C_j$. Conditional on $(D_j,\mathcal H_t)$, both randomizations are independent of $E$. Both agents observe the delivered $H_{t+1}$ and append it to their persistent registers. Future policies read their private record and common transcript through the stated conditional laws; no additional access to old transmission randomization is included.

The sampled $\widetilde Y_{t+1}$ and represented $Y_j$ have the same conditional distribution but are different variables. In particular,
\begin{equation}
 H_{t+1}\mathbin{\perp}(Z,D_{1-j})\mid(D_j,\mathcal H_t).
 \label{imem:conditional-independence}
\end{equation}
A direct sensor $C_j(Y_j,dh)$ changes this conditional independence and therefore changes the experiment.

\propositionheading{Public calibration}{imem:calibration}
For every stage of the declared protocol, the public predictive law is
\begin{equation}
 P(H_{t+1}\in dh\mid\mathcal H_t)
 =\int C_j(y,dh)\Gamma_t^{Y_j}(dy).
 \label{imem:calibration-formula}
\end{equation}
\begin{proof}
Average \cref{imem:message-kernel} over the common marginal of $D_j$. The integral of the conditional law $P(Y_j\in dy\mid D_j,\mathcal H_t)$ is $\Gamma_t^{Y_j}(dy)$. Substitution gives \cref{imem:calibration-formula}.
\end{proof}
This equality concerns the message marginal. Its joint law with $E$ is $\Gamma_t(de)\kappa_{j,t}(d_j(e),dh)$, which generally differs from $\Gamma_t(de)C_j(Y_j(e),dh)$. Equal public output laws need not give equal receiver posteriors.

\section{The history law and exact assimilation}

\theoremheading{Normalized histories and Bayesian recursion}{imem:history-theorem}
For every nonnegative integer $n$, the preparation and \cref{imem:protocol} define the consistent probability law
\begin{equation}
 \mathbb P_n(de,dh_{1:n})
 =\Lambda_0(de)\prod_{t=1}^n
 \kappa_{j_t,t-1}(h_{1:t-1},d_{j_t}(e);dh_t).
 \label{imem:history-law}
\end{equation}
Here the history argument, suppressed in \cref{imem:message-kernel}, specifies the posterior used by that policy. Suppose, for density formulas, that $C_j(y,dh)=c_j(h\mid y)\tau(dh)$ for a fixed sigma-finite $\tau$. Set
\begin{equation}
 k_{j,t}(h\mid d_j)=\int c_j(h\mid y)S_{j,t}(d_j,dy),\qquad
 \bar k_{j,t}(h)=\int k_{j,t}(h\mid d_j(e))\Gamma_t(de).
 \label{imem:message-density}
\end{equation}
At every realized prefix where $0<\bar k_{j,t}(h)<\infty$, the exact common posterior update is
\begin{equation}
 \Gamma_{t+1}(de)
 =\frac{k_{j,t}(h_{t+1}\mid d_j(e))}{\bar k_{j,t}(h_{t+1})}\Gamma_t(de).
 \label{imem:bayes}
\end{equation}
The denominator condition holds for predictive-almost every next message. Without domination, the corresponding update is the disintegration of $\Gamma_t(de)\kappa_{j,t}(d_j(e),dh)$.
\begin{proof}
Start from $\mathbb P_0=\Lambda_0$. At stage $t$, the previously chosen conditional kernels make the next policy a measurable normalized kernel of $(e,h_{1:t})$. Integrating its last message gives one. Induction proves normalization of \cref{imem:history-law} and consistency with each preceding marginal. Conditional on the existing transcript, the joint law of the prepared variables and next message is $\Gamma_t(de)\kappa_{j,t}(d_j(e),dh)$. Its message marginal is \cref{imem:calibration-formula}. Dividing the joint density by \cref{imem:message-density} proves \cref{imem:bayes}. A normalized nonnegative density is finite almost everywhere under its base measure and positive almost surely under its own law. This gives the denominator claim. Disintegrating the same joint measure proves the statement without densities.
\end{proof}

For clarity, private continuous records can also be handled by explicit densities. If $L_i(u,dd_i)=l_i(d_i\mid u)\sigma_i(dd_i)$, with sigma-finite $\sigma_i$, take $\eta=P_0\otimes\sigma_0\otimes\sigma_1$. Then $g_0=d\Gamma_0/d\eta=l_0l_1$, and
\begin{equation}
 g_{t+1}(z,d)=
 \frac{g_t(z,d)k_{j,t}(h_{t+1}\mid d_j)}
 {\int g_t(z',d')k_{j,t}(h_{t+1}\mid d_j')\eta(dz',dd')}.
 \label{imem:dominated-recursion}
\end{equation}
For a private value $d_i$, divide the slice $g_t(z,d_i,d_{1-i})P_0(dz)\sigma_{1-i}(dd_{1-i})$ by its integral. This gives $\Pi_{i,t}$ wherever that integral is positive and finite, without placing a density assumption on $P_0$ itself.

\propositionheading{Sender invariance and the receiver likelihood}{imem:sender-receiver}
If agent $j$ sends at step $t+1$, then $\Pi_{j,t+1}=\Pi_{j,t}$ almost surely. For the receiving agent $i=1-j$, define
\begin{equation}
 R_{ij,t}(dd_j\mid z_i)
 =P(D_j\in dd_j\mid Z_i=z_i,D_i=d_i,\mathcal H_t).
 \label{imem:peer-private-conditional}
\end{equation}
Under the dominated hypotheses of \cref{imem:history-theorem}, its actual local likelihood and update are
\begin{equation}
 \begin{aligned}
 \ell_{i\leftarrow j,t}(h\mid z_i)
 &=\int k_{j,t}(h\mid d_j)R_{ij,t}(dd_j\mid z_i),\\
 \zeta_{i,t+1}(dz_i)
 &=\frac{\ell_{i\leftarrow j,t}(h_{t+1}\mid z_i)\zeta_{i,t}(dz_i)}
 {\int\ell_{i\leftarrow j,t}(h_{t+1}\mid z_i')\zeta_{i,t}(dz_i')}.
 \end{aligned}
 \label{imem:receiver-update}
\end{equation}
\begin{proof}
Condition \cref{imem:bayes} further on $D_j=d_j$. Its message factor is constant in every remaining variable, so it cancels from the conditional numerator and denominator. This proves invariance of the entire private posterior; \cref{imem:conditional-independence} gives the same conclusion without densities. For the receiver, first condition the actual message kernel on $(Z_i,D_i,\mathcal H_t)$. The remaining private coordinate has law \cref{imem:peer-private-conditional}; integrating it gives the first equation of \cref{imem:receiver-update}. Marginalizing the private joint update then gives the second.
\end{proof}

The conditional law $R_{ij,t}$ carries the relationship between the receiver's represented state and the sender's evidence. A local marginal $\zeta_{i,t}$ does not in general specify it. If a message policy reads only common information, its factor is independent of both private values and cancels from the common posterior as well. A deterministic echo of an already registered message has this property in the kernel formulation. A fresh posterior sample can instead convey the speaker's private evidence to its peer, while leaving the speaker's own posterior unchanged.

There is an exact information restriction for the fixed preparation. Every factor in \cref{imem:history-law} depends on $E$ only through $D$ at a fixed transcript. Thus $Z$ and $H_{1:n}$ are conditionally independent given $D$, and, when the mutual informations are finite,
\begin{equation}
 I(Z;H_{1:n}\mid D_i)
 =I(Z;D_{1-i}\mid D_i)
 -I(Z;D_{1-i}\mid D_i,H_{1:n}).
 \label{imem:information-bound}
\end{equation}
Indeed, expand $I(Z;(D_{1-i},H_{1:n})\mid D_i)$ in the two possible orders; the term conditioned on both private records vanishes. The protocol redistributes evidence from the preparation. New measurements require an additional observation mechanism.

\section{A fixed reconstruction and the retained state}

Disintegrate the preparation once as
\begin{equation}
 \Lambda_0(dz,dd)=W(d,dz)\lambda_0(dd),\qquad
 \lambda_0=P(D\in\cdot),\quad W(d,\cdot)=P(Z\in\cdot\mid D=d).
 \label{imem:fixed-reconstruction}
\end{equation}
The null-set version of $W$ is fixed with the model. Write $W_i$ and $W^{Y_j}$ for its local-state and result marginals, and let
\begin{equation}
 V_j(d,dh)=\int C_j(y,dh)W^{Y_j}(d,dy).
 \label{imem:encoded-reconstruction}
\end{equation}

\theoremheading{Exact private-record reduction}{imem:reduction}
For every finite stage of the protocol,
\begin{equation}
 \Gamma_t(dz,dd)=W(d,dz)\lambda_t(dd),\qquad
 \lambda_t=P(D\in\cdot\mid\mathcal H_t).
 \label{imem:reconstructed-posterior}
\end{equation}
For its chosen conditional versions the local recognition and transmission law are
\begin{equation}
 \begin{aligned}
 \zeta_i^{\lambda_t,d_i}(dz_i)
 &=\int W_i(d,dz_i)\lambda_t(dd\mid D_i=d_i),\\
 \kappa_j^{\lambda_t}(d_j,dh)
 &=\int V_j(d,dh)\lambda_t(dd\mid D_j=d_j).
 \end{aligned}
 \label{imem:reduced-readouts}
\end{equation}
If the latter kernel has density $k_j^{\lambda_t}(h\mid d_j)$ relative to $\tau$, the common state obeys the closed recursion
\begin{equation}
 \lambda_{t+1}(dd)=
 \frac{k_j^{\lambda_t}(h_{t+1}\mid d_j)\lambda_t(dd)}
 {\int k_j^{\lambda_t}(h_{t+1}\mid d_j')\lambda_t(dd')}.
 \label{imem:reduced-update}
\end{equation}
Reconstruction by $W$ makes this recursion exactly equivalent to \cref{imem:bayes} for the complete joint posterior, and hence for every finite future sequence of protocol messages.
\begin{proof}
At a fixed history, its likelihood $\mathcal L_t(d)$ is a function of $D$ alone. In the dominated case, for bounded measurable $f$ and $a$ the posterior integral is
\begin{equation}
 \begin{split}
 \int f(z)a(d)\Gamma_t(dz,dd)
 &=\frac{\int a(d)\mathcal L_t(d)
       [\int f(z)W(d,dz)]\lambda_0(dd)}
       {\int\mathcal L_t(d)\lambda_0(dd)}\\
 &=\int a(d)\left[\int f(z)W(d,dz)\right]\lambda_t(dd).
 \end{split}
 \label{imem:reconstruction-proof}
\end{equation}
Products of bounded measurable functions determine the joint law by a monotone-class argument. Without domination, the same identity follows by integrating $f(z)$ first in \cref{imem:history-law}; its conditional history kernel depends on $d$ alone. This proves \cref{imem:reconstructed-posterior}. Conditioning it on the observed private coordinate and taking the indicated marginals gives \cref{imem:reduced-readouts}. Multiplication by the message density changes only $\lambda_t$, proving \cref{imem:reduced-update}. Conversely, marginalizing \cref{imem:bayes} gives that update. At each future step the reconstruction supplies the same message kernel and conditional update, so induction proves equality of every finite future joint law.
\end{proof}

The sufficient state for an agent includes the fixed preparation, reconstruction and encoders, its private value, the joint $\lambda_t$, and the next speaker. When queries of actual past records are admitted, take the retained state to be
\begin{equation}
 \mathcal R_{i,t}
 =\bigl(P_0,L_0,L_1,W,C_0,C_1,d_i,\lambda_t,j_{t+1},H_{1:t}\bigr).
 \label{imem:retention}
\end{equation}
Equality of these states gives the same current private posterior and, by \cref{imem:reduction}, the same joint conditional laws for all future protocol messages and represented-variable readouts. The last coordinate preserves the literal transcript; its values need not be recoverable from $\lambda_t$. This separates inferential sufficiency from persistence of the full admitted record interface.

On the family $\{W\lambda\}$, the map $W\lambda\mapsto\lambda$ is invertible by the fixed reconstruction. It removes redundant changing coordinates, rather than asserting a lossy parent compression in the sense of \cref{inh:two-retentions}. The fixed $W$ can be large, and a law on continuous private-record spaces need not have a finite description. Keeping only the two private-record marginals would also discard the conditional pair law in \cref{imem:reduced-readouts}.

The supplied coupling has a separate role. If the preparation factorizes across $(Z_0,D_0)$ and $(Z_1,D_1)$, its fixed-history likelihood is a product of a function of $D_0$ and a function of $D_1$, even though the factors use the observed prefix. The common posterior consequently remains a product across the two agents. After conditioning on its own private record, an agent's own represented-state posterior cannot change when its peer speaks. It may learn about the peer, but information transfer to its own target requires a cross-agent relation supplied by the model.

\section{Free energies for the declared information}

Fix a dominated observed history $h_{1:n}$, and define
\begin{equation}
 \mathcal L_n(d;h)=\prod_{t=1}^n
 k_{j_t,t-1}(h_t\mid d_{j_t};h_{1:t-1}),\qquad
 \mathcal Z_n(h)=\int\mathcal L_n(d;h)\Lambda_0(dz,dd).
 \label{imem:history-likelihood}
\end{equation}
The policy values at each observed prefix are those determined from the actual preparation and protocol. They remain fixed when a candidate recognition law is varied. Assume $0<\mathcal Z_n(h)<\infty$. For a candidate probability $Q\ll\Lambda_0$ supported where $\mathcal L_n>0$, define
\begin{equation}
 \begin{split}
 \mathcal F_n(Q;h)
 &=\E_Q\left[
 \log\frac{dQ}{d\Lambda_0}(E)-\log\mathcal L_n(D;h)\right]\\
 &=\KL(Q\Vert\Lambda_0)-\E_Q\log\mathcal L_n(D;h).
 \end{split}
 \label{imem:full-free-energy}
\end{equation}
In the statements below, the relative entropy and the absolute expectation of the displayed log likelihood are finite. These conditions prevent an undefined subtraction of infinities and imply the finite posterior relative entropy used in each identity. The statements extend to any other admissible domain on which the defining logarithmic integral and chain rules are well defined, but no such extension is needed here.

\propositionheading{Full, private, and local inference objectives}{imem:free-energy}
Under these conditions,
\begin{equation}
 \mathcal F_n(Q;h)=\KL(Q\Vert\Gamma_n)-\log\mathcal Z_n(h).
 \label{imem:full-free-energy-identity}
\end{equation}
For agent $i$, let $\Lambda_{0,i}^{d_i}=P(E\in\cdot\mid D_i=d_i)$ and
$\mathcal Z_{i,n}=\int\mathcal L_n(D;h)\Lambda_{0,i}^{d_i}(dE)$, assumed positive and finite. Its private objective is
\begin{equation}
 \begin{split}
 \mathcal F_{i,n}(Q;d_i,h)
 &=\E_Q\left[
 \log\frac{dQ}{d\Lambda_{0,i}^{d_i}}(E)-\log\mathcal L_n(D;h)\right]\\
 &=\KL(Q\Vert\Pi_{i,n})-\log\mathcal Z_{i,n}.
 \end{split}
 \label{imem:private-free-energy}
\end{equation}
Let $\zeta_{i,0}$ denote its initial local marginal and put
\begin{equation}
 b_{i,n}(z_i)=\E_{\Lambda_{0,i}^{d_i}}
 [\mathcal L_n(D;h)\mid Z_i=z_i].
 \label{imem:local-history-likelihood}
\end{equation}
For a candidate local correlated law $\mathcal Q_i(da,dk)=s_i(da)q_i(dk\mid a)$, where $k=(u,v)$, define
\begin{equation}
 \mathcal F_{i,n}^{\mathrm{loc}}(\mathcal Q_i)
 =\E_{\mathcal Q_i}\left[
 \log\frac{d\mathcal Q_i}{d\zeta_{i,0}}(Z_i)-\log b_{i,n}(Z_i)\right].
 \label{imem:local-free-energy}
\end{equation}
With the same absolute-continuity and integrability hypotheses on this local space,
\begin{equation}
 \begin{split}
 \mathcal F_{i,n}^{\mathrm{loc}}(\mathcal Q_i)
 ={}&\KL(s_i\Vert s_{i,n})\\
 &+\int\KL(q_i(\cdot\mid a)\Vert q_{i,n}(\cdot\mid a))s_i(da)
 -\log\mathcal Z_{i,n}.
 \end{split}
 \label{imem:local-free-energy-split}
\end{equation}
Whenever the corresponding posterior satisfies the stated finiteness conditions, each unrestricted minimum is its negative log evidence, attained only by that posterior up to null sets.
\begin{proof}
Bayes' formula gives $d\Gamma_n/d\Lambda_0=\mathcal L_n/\mathcal Z_n$. Subtract its logarithm from $\log(dQ/d\Lambda_0)$ inside \cref{imem:full-free-energy}; integrating yields \cref{imem:full-free-energy-identity}. The same argument on the conditional base $\Lambda_{0,i}^{d_i}$ proves \cref{imem:private-free-energy}. Marginalization gives
$\zeta_{i,n}(dz_i)=b_{i,n}(z_i)\zeta_{i,0}(dz_i)/\mathcal Z_{i,n}$, so the local objective equals $\KL(\mathcal Q_i\Vert\zeta_{i,n})-\log\mathcal Z_{i,n}$. To split that relative entropy, disintegrate both laws over $A_i$. Their logarithmic density ratio is the marginal label ratio plus the conditional content ratio. Integration gives \cref{imem:local-free-energy-split}. Relative entropy is nonnegative and vanishes precisely for equal probability laws, proving the constrained minima over normalized laws. When the posterior belongs to the stated finite-objective domain, these minima are attained there.
\end{proof}

The last integrability qualification matters for general laws. Finiteness of the evidence alone does not imply finite expected log likelihood at the posterior. The normalized posterior and its relative-entropy characterization still exist under the evidence condition; the separated finite expectations in \cref{imem:full-free-energy} require the additional hypothesis stated above.

The conditional information omitted by a local objective is explicit. For an admissible full $Q$, disintegration over $Z_i$ gives
\begin{equation}
 \begin{split}
 \KL(Q\Vert\Pi_{i,n})
 ={}&\KL(Q^{Z_i}\Vert\zeta_{i,n})\\
 &+\int\KL\bigl(Q(\cdot\mid Z_i=z_i)
 \Vert\Pi_{i,n}(\cdot\mid Z_i=z_i)\bigr)Q^{Z_i}(dz_i).
 \end{split}
 \label{imem:omitted-conditional}
\end{equation}
The proof is the same logarithmic split. Optimizing only the local term does not identify the conditional law of the peer's private record required in \cref{imem:receiver-update}. Likewise, a product restriction on presentation and content changes the admitted variational family; it is not part of the exact posterior construction.

\propositionheading{Exact elimination in free energy}{imem:reduced-free-energy}
Let $Q(dz,dd)=Q^Z(d,dz)\nu(dd)$, with the absolute-continuity and finiteness conditions of \cref{imem:free-energy}. Then
\begin{equation}
 \begin{split}
 \mathcal F_n(Q;h)
 ={}&\KL(\nu\Vert\lambda_n)
 +\int\KL(Q^Z(d,\cdot)\Vert W(d,\cdot))\nu(dd)
 -\log\mathcal Z_n(h).
 \end{split}
 \label{imem:conditional-free-energy}
\end{equation}
For fixed $\nu$, the minimizing conditional is $W$ almost everywhere, whenever that candidate is admissible. The remaining objective is exactly
\begin{equation}
 \begin{split}
 \mathcal F_n^D(\nu;h)
 &=\E_\nu\left[
 \log\frac{d\nu}{d\lambda_0}(D)-\log\mathcal L_n(D;h)\right]\\
 &=\KL(\nu\Vert\lambda_n)-\log\mathcal Z_n(h).
 \end{split}
 \label{imem:private-record-free-energy}
\end{equation}
\begin{proof}
Substitute $\Gamma_n=W\lambda_n$ into \cref{imem:full-free-energy-identity}. Splitting its logarithmic ratio over $D$ gives the marginal relative entropy and the averaged conditional relative entropy in \cref{imem:conditional-free-energy}. The latter is minimized uniquely by $Q^Z=W$, $\nu$-almost everywhere. Because $\mathcal L_n$ depends only on $D$, the defining expectation for $Q=W\nu$ reduces to the first line of \cref{imem:private-record-free-energy}. Bayes' formula for $\lambda_n$ gives its second line.
\end{proof}

These are unit-temperature inference objectives for a fixed observed history. Their evidence terms are constant in the candidate law within that optimization block. The agents condition on different private information, so their objectives have different minimizers. Recomputing historical policies from the candidate $Q$ inside \cref{imem:history-likelihood} would also vary the observation model and define a different optimization problem. No natural-gradient update, monotone sum during communication, or bank-learning rule follows from the variational identities alone.

\section{An exact two-way binary exchange}

Let every variable be a bit. Draw $X_0$ fairly, set $P(X_1=X_0\mid X_0)=3/4$, and, independently conditional on the cues, prepare
\begin{equation}
 P(A_i=X_i\mid X_i)=\frac34,\qquad
 P(Y_i=1-X_i\mid X_i)=\frac34,\qquad
 P(D_i=X_i\mid X_i)=\frac34.
 \label{imem:binary-preparation}
\end{equation}
All base measures are counting measures, and all preparation atoms are positive. Since each cue is fair, $A_i$ is fair and $P(X_i=A_i\mid A_i)=3/4$. This is a local bank of the form \cref{imem:bank}, with a negative cue/result response and correlated presentation/content variables. Use identity encoders on the sampled result bits, with agent 1 speaking first.

At any information state of this protocol, if $p_i=P(X_i=1\mid\mathcal I_i)$, then
\begin{equation}
 \begin{aligned}
 P(A_i=1\mid\mathcal I_i)&=\frac14+\frac12p_i,&
 P(Y_i=1\mid\mathcal I_i)&=\frac34-\frac12p_i,\\
 \operatorname{Cov}(A_i,X_i\mid\mathcal I_i)&=\frac12p_i(1-p_i),&
 \operatorname{Cov}(Y_i,X_i\mid\mathcal I_i)&=-\frac12p_i(1-p_i).
 \end{aligned}
 \label{imem:binary-correlations}
\end{equation}
Indeed, conditioning further on the cues leaves the independent presentation and result kernels in \cref{imem:binary-preparation} unchanged: the message likelihood depends on the preparation only through the private records. Total expectation and covariance then give \cref{imem:binary-correlations}. In particular, the presentation/content covariance is nonzero at every finite history.

Initially the speaker's result-sampling probabilities are $k_{1,0}(1\mid0)=5/8$ and $k_{1,0}(1\mid1)=3/8$. Integrating its private sensor and then the cue coupling gives
\begin{equation}
 \begin{aligned}
 P(H_1=1\mid X_1=0)&=\frac9{16},&
 P(H_1=1\mid X_1=1)&=\frac7{16},\\
 P(H_1=1\mid X_0=0)&=\frac{17}{32},&
 P(H_1=1\mid X_0=1)&=\frac{15}{32}.
 \end{aligned}
 \label{imem:binary-first-likelihood}
\end{equation}
For example, the first probability is $(3/4)(5/8)+(1/4)(3/8)$, and the third is $(3/4)(9/16)+(1/4)(7/16)$. At $D_0=1,H_1=1$, Bayes' rule gives
\begin{equation}
 P(X_0=1\mid D_0=1,H_1=1)
 =\frac{(3/4)(15/32)}{(3/4)(15/32)+(1/4)(17/32)}
 =\frac{45}{62}.
 \label{imem:binary-first-posterior}
\end{equation}
At $D_0=0$ the same calculation gives $5/22$. The reply therefore uses
\begin{equation}
 k_{0,1}(1\mid D_0=0,H_1=1)=\frac7{11},\qquad
 k_{0,1}(1\mid D_0=1,H_1=1)=\frac{12}{31}.
 \label{imem:binary-reply}
\end{equation}
Both values differ from the initial sampling probabilities. The reception changes the outgoing law.

The four private-record weights make the second reception explicit. In the order $(00,01,10,11)$, direct marginalization of \cref{imem:binary-preparation}, followed by the two reweightings, yields
\begin{equation}
 \begin{aligned}
 \lambda_0&=\frac1{32}(9,7,7,9),\\
 \lambda_1\bigm|_{H_1=1}&=\frac1{128}(45,21,35,27),\\
 \lambda_2\bigm|_{H_1=H_2=1}
 &=\left(\frac{105}{242},\frac{49}{242},\frac{70}{341},\frac{54}{341}\right).
 \end{aligned}
 \label{imem:binary-private-weights}
\end{equation}
The first normalization is $P(H_1=1)=1/2$. The second is
$(66/128)(7/11)+(62/128)(12/31)=33/64$; multiplying the corresponding rows of $\lambda_1$ by \cref{imem:binary-reply} and dividing by this value proves the last line. The fixed reconstruction supplies
\begin{equation}
 \begin{aligned}
 P(X_0=1\mid D)&=\left(\frac16,\frac5{14},\frac9{14},\frac56\right),\\
 V_0(D;\{1\})&=\left(\frac23,\frac47,\frac37,\frac13\right),\\
 V_1(D;\{1\})&=\left(\frac23,\frac37,\frac47,\frac13\right).
 \end{aligned}
 \label{imem:binary-reconstruction}
\end{equation}
For the first table, each entry is the normalized sum over $X_1$ of the cue prior and two private-sensor factors. The corresponding table for $X_1$ exchanges $D_0,D_1$; the last two tables then follow from $P(Y_i=1\mid X_i)=3/4-X_i/2$.

Take actual private values $D_0=1,D_1=0$ and delivered bits $H_1=H_2=1$. Each agent still conditions only on its own private value. The common weights give
\begin{equation}
 P(D_0=1\mid D_1=0,H_1=H_2=1)=\frac{308}{959},
\end{equation}
and hence agent 1's final cue probability is
\begin{equation}
 P(X_1=1\mid D_1=0,H_1=H_2=1)
 =\frac{651}{959}\frac16+\frac{308}{959}\frac5{14}
 =\frac{437}{1918}.
 \label{imem:binary-second-posterior}
\end{equation}
The accessible cue probabilities are therefore
\begin{equation}
 \begin{array}{c|cc}
 \text{Registered messages}&\text{Agent 0, }D_0=1&\text{Agent 1, }D_1=0\\ \hline
 \text{None}&3/4&1/4\\
 H_1=1&45/62&1/4\\
 H_1=H_2=1&45/62&437/1918.
 \end{array}
 \label{imem:binary-exchange}
\end{equation}
The sender invariance in each row follows from \cref{imem:sender-receiver}; the receiver changes at both receptions. Presentation and result probabilities follow from \cref{imem:binary-correlations}. The fixed reconstruction and four changing masses reproduce the full joint law, including its local correlations.

The same preparation distinguishes two errors in message interpretation. A direct observation of represented $Y_1=1$ has cue likelihoods $5/8$ and $3/8$ conditional on $X_0=0,1$, respectively. It gives $P(X_0=1\mid D_0=1,Y_1=1)=9/14$, in place of $45/62$. If agent 1 instead sends two fresh samples before a reply, its private posterior is unchanged, so the samples are independent conditional on $D_1$. Their probabilities conditional only on $X_1$ are
\begin{equation}
 P(H_1=H_2=1\mid X_1=0)=\frac{21}{64},\qquad
 P(H_1=H_2=1\mid X_1=1)=\frac{13}{64}.
 \label{imem:binary-repeated-samples}
\end{equation}
These equal $(3/4)(5/8)^2+(1/4)(3/8)^2$ and its exchanged version. Integrating the cue coupling and conditioning on $D_0=1$ gives $45/64$. Squaring the once-marginalized likelihoods $9/16,7/16$ instead gives $171/244$, because that operation deletes their shared private cause. A deterministic echo $H_2=H_1$ has conditional likelihood one at the registered bit and causes no further update.

\section{Gaussian communication and an exact affine closure}

Take finite presentation sets, vector cues $X_i\in\R^{d_i}$ and results $Y_i\in\R^{e_i}$, and a strictly positive label table $p_{ab}$ of total mass one, with marginals $\pi_{0a},\pi_{1b}$. For each pair choose the conditional cue law
\begin{equation}
 \begin{pmatrix}X_0\\X_1\end{pmatrix}\Bigm|A_0=a,A_1=b
 \sim\mathcal N\left(
 \begin{pmatrix}\nu_{0a}\\\nu_{1b}\end{pmatrix},
 \begin{pmatrix}L_{0a}&S_{ab}\\S_{ab}^T&L_{1b}\end{pmatrix}\right),
 \label{imem:gaussian-coupling}
\end{equation}
where the block covariance is positive definite. For instance,
$S_{ab}=L_{0a}^{1/2}J_{ab}L_{1b}^{1/2}$ with $\|J_{ab}\|<1$ satisfies this condition. Given the cues and labels, generate the results independently by
\begin{equation}
 Y_i\mid X_i=u,A_i=a
 \sim\mathcal N\bigl(\omega_{ia}+F_{ia}(u-\nu_{ia}),Q_{ia}\bigr),
 \qquad Q_{ia}\succ0.
 \label{imem:gaussian-results}
\end{equation}
The matrices $F_{ia}$ are arbitrary. Summing the label table and integrating the normalized cue and result kernels proves normalization. The marginal cue law for agent $i$ does not depend on the other label, so its local bank is precisely
$\pi_{ia}\mathcal N(du;\nu_{ia},L_{ia})K_{ia}(u,dv)$. The label coupling and cross-covariances are additional model data.

Let the private sensors be $D_i=X_i+\epsilon_i$, where $\epsilon_i\sim\mathcal N(0,N_i)$, $N_i\succ0$, are mutually independent and independent of the prepared variables. This notation distinguishes sensor covariance from the encoder kernel $C_i$. Before any message, the sender, agent 1, has
\begin{equation}
 \begin{aligned}
 T_b&=L_{1b}+N_1,&
 w_b(d)&=\frac{\pi_{1b}\mathcal N(d;\nu_{1b},T_b)}
 {\sum_c\pi_{1c}\mathcal N(d;\nu_{1c},T_c)},\\
 G_b&=L_{1b}T_b^{-1},&
 P_b&=L_{1b}-L_{1b}T_b^{-1}L_{1b},\\
 m_b(d)&=\omega_{1b}+F_{1b}G_b(d-\nu_{1b}),&
 \Sigma_b^Y&=Q_{1b}+F_{1b}P_bF_{1b}^T.
 \end{aligned}
 \label{imem:gaussian-private-posterior}
\end{equation}
Completing the cue quadratic gives $P_b=(L_{1b}^{-1}+N_1^{-1})^{-1}\succ0$. Applying the affine result kernel gives the mixture with weights $w_b$, means $m_b$, and covariances $\Sigma_b^Y\succ0$, as in \cref{cmem:gaussian-posterior}.

Use the encoder $H=B\widetilde Y_1+\eta$ with independent $\eta\sim\mathcal N(0,R)$, $R\succ0$. Its actual first-message density conditional on the sender's private record is
\begin{equation}
 k(h\mid d)=\sum_b w_b(d)
 \mathcal N(h;Bm_b(d),R+B\Sigma_b^YB^T).
 \label{imem:gaussian-message-density}
\end{equation}
The sample and encoder noises are independent of the prepared variables conditional on $D_1$. Conditional on a complete represented value $z$, the likelihood is therefore
\begin{equation}
 P(H\in dh\mid Z=z)/dh
 =\int\mathcal N(d;u_1,N_1)k(h\mid d)dd.
 \label{imem:gaussian-integrated-likelihood}
\end{equation}
At private value $d_0$ and received $h$, the augmented posterior is
\begin{equation}
 \Pi(dz,dd_1\mid d_0,h)
 =\frac{P_0(dz)\mathcal N(d_0;u_0,N_0)
       \mathcal N(d_1;u_1,N_1)k(h\mid d_1)}
       {\mathcal Z(d_0,h)}dd_1.
 \label{imem:gaussian-augmented-posterior}
\end{equation}
The normalizer is positive and finite for every finite $(d_0,h)$. Positivity follows from the strictly positive Gaussian factors. Finiteness follows because the private sensor density is bounded and the finite mixture in \cref{imem:gaussian-message-density} is bounded by the largest of its component density maxima, independently of $d$. All integrands are nonnegative, so their order can be exchanged. Integrating $h$, then the private records and prepared variables, gives total mass one.

\propositionheading{The first receiver likelihood for Gaussian banks}{imem:gaussian-receiver}
For the preparation above, put
\begin{equation}
 \begin{aligned}
 u_{ab}(x)&=\nu_{1b}+S_{ab}^TL_{0a}^{-1}(x-\nu_{0a}),\\
 U_{ab}&=N_1+L_{1b}-S_{ab}^TL_{0a}^{-1}S_{ab},\\
 \ell_0(h\mid a,x)&=\sum_b\frac{p_{ab}}{\pi_{0a}}
 \int\mathcal N(d;u_{ab}(x),U_{ab})k(h\mid d)dd.
 \end{aligned}
 \label{imem:gaussian-local-likelihood}
\end{equation}
Then the exact correlated local posterior is
\begin{equation}
 \zeta_0(a,dx,dy\mid d_0,h)
 \propto\pi_{0a}\mathcal N(dx;\nu_{0a},L_{0a})K_{0a}(x,dy)
 \mathcal N(d_0;x,N_0)\ell_0(h\mid a,x).
 \label{imem:gaussian-local-posterior}
\end{equation}
\begin{proof}
Completing the joint cue quadratic conditional on $(A_0,A_1,X_0)=(a,b,x)$ gives the mean $u_{ab}(x)$ and covariance $L_{1b}-S_{ab}^TL_{0a}^{-1}S_{ab}$ for $X_1$. Adding the private sensor noise gives $U_{ab}\succ0$; the Schur complement is positive by \cref{imem:gaussian-coupling}. Since the marginal law of $X_0$ at label $a$ does not depend on $b$, conditioning on $X_0=x$ leaves the label weights $p_{ab}/\pi_{0a}$ unchanged. Integrating the sender's private record against \cref{imem:gaussian-message-density} therefore gives \cref{imem:gaussian-local-likelihood}. The own-private likelihood and this message likelihood are conditionally independent of $Y_0$ at fixed $(A_0,X_0)$. Multiplying by the local bank and normalizing proves \cref{imem:gaussian-local-posterior}, equivalently by marginalizing \cref{imem:gaussian-augmented-posterior}.
\end{proof}

The integral includes the sender's posterior-dependent weights $w_b(d)$. It cannot generally be replaced by a Gaussian observation of $Y_1$ or by evaluating those weights at a message mean.

The scalar, one-presentation case makes the different observation mechanism explicit. Write $\operatorname{Cov}(X_0,X_1)=S$, with $S^2<L_0L_1$, result slope $f_1$, and encoder slope $b$. Let
\begin{equation}
 g=\frac{L_1}{L_1+N_1},\qquad
 P=\frac{L_1N_1}{L_1+N_1},\qquad a=bf_1g.
 \label{imem:scalar-gain}
\end{equation}
A fresh posterior sample has the realization
$\widetilde Y_1=\omega_1+f_1g(D_1-\nu_1)+\xi$, where $\xi$ is independent Gaussian with variance $Q_1+f_1^2P$. Substituting $D_1=X_1+\epsilon_1$ gives
\begin{equation}
 \begin{aligned}
 H&=b\omega_1+a(X_1-\nu_1)+a\epsilon_1+b\xi+\eta,\\
 N_{\mathrm{eff}}&=a^2N_1+b^2(Q_1+f_1^2P)+R,\\
 H\mid X_1=x&\sim
 \mathcal N\bigl(b\omega_1+a(x-\nu_1),N_{\mathrm{eff}}\bigr).
 \end{aligned}
 \label{imem:scalar-message}
\end{equation}
The direct sensor $H_{\mathrm{dir}}=bY_1+\eta$ instead has gain $bf_1$ and conditional variance $b^2Q_1+R$. Nevertheless,
\begin{equation}
 \begin{aligned}
 \operatorname{Var}(H)&=a^2L_1+N_{\mathrm{eff}}
 =b^2(f_1^2L_1+Q_1)+R
 =\operatorname{Var}(H_{\mathrm{dir}}),\\
 \operatorname{Cov}(X_0,H)&=aS,&
 \operatorname{Cov}(X_0,H_{\mathrm{dir}})&=bf_1S.
 \end{aligned}
 \label{imem:scalar-calibration}
\end{equation}
The variance equality uses $g^2(L_1+N_1)+P=L_1$; both means equal $b\omega_1$. Thus even equal Gaussian message marginals leave different cross-agent evidence.

For completeness, first conditioning on $D_0=d_0$ gives
\begin{equation}
 \begin{aligned}
 m_0&=\nu_0+\frac{L_0}{L_0+N_0}(d_0-\nu_0),&
 v_0&=\frac{L_0N_0}{L_0+N_0},\\
 m_1^0&=\nu_1+\frac{S}{L_0+N_0}(d_0-\nu_0),&
 v_1^0&=L_1-\frac{S^2}{L_0+N_0},\\
 c_{01}^0&=\frac{SN_0}{L_0+N_0}.&&
 \end{aligned}
\end{equation}
Let $m_H=b\omega_1+a(m_1^0-\nu_1)$ and $V_H=a^2v_1^0+N_{\mathrm{eff}}$. Conditioning the remaining Gaussian on $H=h$ gives
\begin{equation}
 \begin{aligned}
 m_0^+&=m_0+\frac{ac_{01}^0}{V_H}(h-m_H),&
 v_0^+&=v_0-\frac{a^2(c_{01}^0)^2}{V_H},\\
 Y_0\mid d_0,h&\sim
 \mathcal N\bigl(\omega_0+f_0(m_0^+-\nu_0),Q_0+f_0^2v_0^+\bigr).
 \end{aligned}
 \label{imem:scalar-receiver}
\end{equation}
Both conditioning operations are completions of positive Gaussian quadratics. Positive sensor and encoder noise make their Schur complements positive. These formulas therefore give exact conditioning without a moment projection.

\theoremheading{Repeated affine Gaussian closure}{imem:gaussian-closure}
Suppose each agent has one presentation, the joint cue covariance is positive definite, and all result kernels, private sensors, and encoders are affine Gaussian with positive-definite noise covariances. For every finite protocol history, the common private-record posterior $\lambda_t$ is a nondegenerate joint Gaussian. Its mean and full covariance, together with the fixed Gaussian reconstruction $W$, determine every posterior and future message law in \cref{imem:reduction}.
\begin{proof}
After the fixed singleton labels are suppressed, the preparation makes $(Z,D)$ jointly Gaussian and $\lambda_0$ nondegenerate. Completing its quadratic gives $W(d,dz)$ an affine mean in $d$ and a covariance independent of $d$. Suppose $\lambda_t=\mathcal N(m_t,S_t)$ with $S_t\succ0$. Then $W\lambda_t$ is jointly Gaussian. Conditioning its $Y_j$ coordinate on $D_j=d_j$ gives an affine mean and constant covariance. After fresh sampling and encoding, the message therefore has a representation
\begin{equation}
 H_{t+1}=b_t+A_tD+\eta_t,\qquad
 \eta_t\sim\mathcal N(0,R_t),\quad R_t\succ0,
 \label{imem:gaussian-reduced-observation}
\end{equation}
where $A_t$ reads only the speaker coordinate and $\eta_t$ is independent of $D$ at the fixed prefix. The coefficients are fixed by the current common state. Multiplying the Gaussian prior density by this likelihood completes a quadratic with
\begin{equation}
 \begin{aligned}
 S_{t+1}^{-1}&=S_t^{-1}+A_t^TR_t^{-1}A_t,\\
 m_{t+1}&=S_{t+1}\left[S_t^{-1}m_t+A_t^TR_t^{-1}(h_{t+1}-b_t)\right].
 \end{aligned}
 \label{imem:gaussian-reduced-recursion}
\end{equation}
The precision remains positive definite. Induction gives the claimed closure for every finite stage; \cref{imem:reduction} gives the reconstruction and future-law assertion.
\end{proof}

The retained covariance in this theorem is the joint covariance of $(D_0,D_1)$. Replacing it by separate-agent blocks changes the conditional averages in \cref{imem:reduced-readouts}. The Gaussian family is an exact special case of the general filter, rather than an assumed closure for multiple presentations.

\section{Failure of finite-mixture closure}

\propositionheading{A two-presentation nonclosure example}{imem:mixture-nonclosure}
There is a preparation of the form \cref{imem:gaussian-coupling,imem:gaussian-results}, with one receiver presentation, two sender presentations, and strictly positive Gaussian noises, whose receiver cue posterior after one message is not any finite positive mixture of nondegenerate Gaussian laws.
\begin{proof}
Let the sender label be $b\in\{-1,1\}$ with probabilities $1/2$, receiver cue mean zero, sender cue means $b\alpha$, $\alpha>0$, common variances $L_0,L_1>0$, and cross-covariance $S>0$ with $S^2<L_0L_1$. Choose sender result means $bm$, $m>0$, common result variance $Q>0$, and slope $F_{1b}=0$. Use the identity encoder with noise variance $R>0$, and positive private-sensor variances $N_0,N_1$. The sender's positive-label posterior weight is
\begin{equation}
 w_+(d)=\frac1{1+\exp[-2\alpha d/(L_1+N_1)]}.
 \label{imem:nonclosure-logistic}
\end{equation}
This follows by taking the ratio of the equal-variance Gaussian label evidences. For any received $h>0$, let $c_\pm=\mathcal N(h;\pm m,Q+R)$. Both constants are positive and $c_+>c_-$, so
\begin{equation}
 k(h\mid d)=c_-+(c_+-c_-)w_+(d).
 \label{imem:nonclosure-message}
\end{equation}
Conditional on $X_0=x$ and $b$, the sender-private record has mean $b\alpha+(S/L_0)x$ and variance $L_1-S^2/L_0+N_1>0$. The label weights remain $1/2$ because the receiver cue marginal is identical under both labels. Represent this conditional record as its moving mean plus a fixed centered Gaussian. As $x\to+\infty$ or $x\to-\infty$, the bounded function $w_+$ tends pointwise to one or zero. Bounded convergence inside \cref{imem:gaussian-local-likelihood} gives
\begin{equation}
 \lim_{x\to+\infty}\ell_0(h\mid x)=c_+,
 \qquad
 \lim_{x\to-\infty}\ell_0(h\mid x)=c_-.
 \label{imem:nonclosure-tails}
\end{equation}
Before the message, $X_0\mid D_0=d_0$ has a nondegenerate Gaussian density $g_0$. Afterwards its density is
\begin{equation}
 p^+(x)=\frac{g_0(x)\ell_0(h\mid x)}{Z_h},\qquad
 Z_h=\int g_0(x)\ell_0(h\mid x)dx.
 \label{imem:nonclosure-posterior}
\end{equation}
The likelihood lies between $c_-$ and $c_+$, so $Z_h$ is positive and finite. Consequently $p^+/g_0$ has two distinct positive finite tail limits.

Suppose instead $p^+=\sum_{r=1}^M a_r\mathcal N(\mu_r,v_r)$, with $a_r>0$ and $v_r>0$, and write $g_0=\mathcal N(\mu_0,v_0)$. A component with $v_r>v_0$ has density ratio to $g_0$ diverging at both tails, because its logarithm has positive quadratic coefficient. A component with $v_r=v_0$ and $\mu_r\ne\mu_0$ diverges at one tail, because its logarithmic ratio is nonconstant affine. Positive mixture weights prevent cancellation, so neither type is compatible with finite limits. Every component with $v_r<v_0$ has ratio tending to zero at both tails. The only remaining components with nonzero finite limits are identical to $g_0$, and their total weight contributes the same limit at both tails. This contradicts \cref{imem:nonclosure-tails}. Zero-weight components can be removed, so the argument covers every finite Gaussian mixture.
\end{proof}

This counterexample concerns closure in a specified finite-mixture family. The exact normalized posterior remains \cref{imem:nonclosure-posterior}; its integral representation and the general private-record filter remain valid. Gaussian prior components, sensors, and encoders alone do not imply a finite Gaussian-mixture posterior when the transmission policy contains posterior-dependent mixing weights.

\section{Local beliefs do not generally determine the next update}

\propositionheading{A reachable collision of local-only retention}{imem:local-retention-failure}
There is one fixed finite preparation, a public preparation record, and a posterior-sampling policy with two positive-probability registered histories having the same pair of complete local recognition laws and private values, but different receiver updates under the same next message. Thus retaining the local joint laws while deleting the distinguishing public record and cross-agent dependence does not define an exact closed assimilation rule.
\begin{proof}
Draw $X_0$ fairly and an independent bit $B$ with $P(B=0)=\rho$, $0<\rho<1$, and set $X_1=X_0\mathbin{\oplus}B$. Conditional on the cues, generate presentations and results as in \cref{imem:binary-preparation}. Each local bank is unchanged. Register the deterministic public preparation readout
\begin{equation}
 C=X_0\mathbin{\oplus}X_1=B.
 \label{imem:parity-record}
\end{equation}
Let $D_0=\ast$ be uninformative and let $P(D_1=X_1\mid X_1)=3/4$, with independent sensor noise. Both parity records occur with positive probability. Given either parity, each individual cue remains fair. The local laws of $(A_i,X_i,Y_i)$ are consequently identical under $C=0$ and $C=1$. The entire local law of $(A_1,X_1,Y_1,D_1)$ is independent of $C$, so the sender's local posterior at either fixed value $D_1=d$ is also identical. Agent 0's local posterior and the actual private values are the same in these two histories as well.

Agent 1 sends a fresh posterior sample $H$ of $Y_1$. Its probabilities are
\begin{equation}
 P(H=1\mid D_1=0,C)=\frac58,\qquad
 P(H=1\mid D_1=1,C)=\frac38.
 \label{imem:parity-message}
\end{equation}
Integrating the private sensor gives $P(H=1\mid X_1=0,C)=9/16$ and $P(H=1\mid X_1=1,C)=7/16$. Thus $P(H=1\mid C)=1/2$ for either parity. At the same received $h=1$, however,
\begin{equation}
 P(X_0=1\mid C=0,H=1)=\frac7{16},\qquad
 P(X_0=1\mid C=1,H=1)=\frac9{16}.
 \label{imem:parity-collision}
\end{equation}
For $C=0$ the cues agree; for $C=1$ they are opposite. Applying \cref{imem:binary-correlations} yields distinct presentation and result marginals,
\begin{equation}
 \begin{array}{c|cc}
 &C=0&C=1\\ \hline
 P(A_0=1\mid C,H=1)&15/32&17/32\\
 P(Y_0=1\mid C,H=1)&17/32&15/32.
 \end{array}
 \label{imem:parity-local-output}
\end{equation}
Let $\mathcal T$ retain the pair of local joint recognition laws, fixed local banks, private values, and schedule, while deleting $C$ and the corresponding cross-agent dependence. The two histories have equal $\mathcal T$. Their exact updates after the same admissible message have different $\mathcal T$ by \cref{imem:parity-collision,imem:parity-local-output}. They therefore violate the retained-fiber condition of \cref{inh:deterministic-descent}, which is necessary for an update defined solely on $\mathcal T$.
\end{proof}

Keeping $C$ resolves this particular collision. In the private-record representation, one may first condition the preparation on $C$ and retain $W_C=P(Z\in\cdot\mid D,C)$ together with $P(D\in\cdot\mid C,H)$. Alternatively, one may include the preparation record among the retained coordinates. A complete common posterior also preserves its effect. The counterexample proves neither minimality of that posterior nor failure for every fixed initialization; it gives two reachable histories of the specified preparation. It uses exact local label/content laws throughout, so its lost information is cross-agent and historical even before any moment projection.

For comparison, if a local label/result posterior $\Pi_{A,Y}$ has label marginal $w$, finite second moments with nondegenerate covariance, and Gaussian result moment match $q_G$, then, whenever the relative entropies are finite,
\begin{equation}
 \KL(\Pi_{A,Y}\Vert w\otimes q_G)
 =I_\Pi(A;Y)+\KL(\Pi_Y\Vert q_G).
 \label{imem:projection-loss}
\end{equation}
Insert $w\otimes\Pi_Y$ into the logarithmic density ratio and integrate to obtain the identity. This separates local label/result dependence from non-Gaussian result information, as in \cref{cmem:projection}. It does not supply the additional private-evidence conditional in \cref{imem:peer-private-conditional}. An approximation to the common state must specify that reconstruction as well. A Gaussian output stored alongside a complete reconstructible transcript is a different retention from one that deletes the distinguishing information.

\openproblemheading{Smaller sufficient states and changing preparations}{imem:open}
For multiple presentations, determine a smaller protocol-specific state that preserves every admitted future message law and receiver update, or specify an approximation together with the conditional reconstruction it changes. The exact joint-private-record representation does not establish minimality or finite descriptive complexity on continuous spaces. New private observations, adaptation of the shared model, differing agent models, or unknown transmission policies require their own normalized mechanisms. These extensions, and any inherited recursive formation, are separate from the fixed-preparation communication theorem proved here.
